\project{10}{PPK2 energy characterisation bench}{Raspberry Pi 4 + Nordic PPK2 + ESP32 NodeMCU as the DUT}{microamp figures}

\begin{unique}
Owns microamp figures. No telemetry product.
\end{unique}

\begin{blueprint}
Draft P9 put the PPK2 in series with a SIM7070 HAT that was still powered from the Pi 5 V
pins. That measures nothing useful. Either lift the HAT 5 V pins and feed VOUT into the
modem power input, or measure a flying-lead DUT (ESP32) that is not also USB-powered. PPK2
source mode is 5 V-class and about 1 A; an LTE TX peak can exceed that, so do not source a
Cat-4 power amplifier from the PPK2.
\end{blueprint}

\subsection*{Intent}

Three firmware states on the ESP32: radio-off idle, BLE advertise, and deep-sleep. The
Pi 4 logs the PPK2 stream and you write down two numbers per state, a median and a p95.
That is the product of this lab. There is no telemetry here, no broker and no uplink: the
deliverable is a small table of current figures that somebody else can act on.

The single idea that makes the measurement valid is that the DUT has exactly one source of
power, and the PPK2 is it. An ESP32 that is also plugged into USB draws most of its current
through the USB rail, and the PPK2 then reports the small remainder that happens to come
through its own output. Traces taken that way look reasonable and mean nothing. Pulling the
DUT USB cable is not a detail of the setup; it is the experiment.

\begin{kit}
Pi 4, PPK2, ESP32 NodeMCU as the DUT, POW-BB, jumpers.
\end{kit}

\subsection*{System architecture}

\diagram{p10_arch}{The measurement path and the control path. The PPK2 sits in the DUT
power rail as a source-measure instrument, and it reports samples to the Pi over its own
USB cable. The DUT's own USB is not part of the picture, because it is not connected.}

Read the figure as two separate cables that meet only inside the PPK2. The power path runs
from a bench 5 V supply into PPK2 VIN, through the instrument, and out of VOUT into the
ESP32 3V3 pin, in source mode, set to 3300 mV. The measurement path runs from the PPK2 over
USB to the Pi 4, which logs samples at 100 kS/s. Ground is common across the PPK2, the DUT
and the Pi, because a current measurement with two ground references is not a measurement.

The instrument has limits and they are part of the design. The PPK2 in source mode is
5 V-class and about 1 A. That is ample for an ESP32 in any of the three states, and it is
not ample for a cellular power amplifier: a Cat-4 transmit peak can exceed it. This is why
the lab measures a flying-lead ESP32 and not the modem from P01, and why the source's
correction of the blueprint draft is repeated on the schematic itself.

\subsection*{Wiring and schematic}

\diagram{p10_schematic}{The whole lab in one drawing. Bench 5 V into VIN, VOUT to the ESP32
3V3 pin at 3300 mV in source mode, one common ground, PPK2 USB to the Pi, and the ESP32 USB
cable drawn where it is not: disconnected.}

\begin{table}[H]\centering\small
\begin{tabular}{p{34mm}p{34mm}p{26mm}p{56mm}}
\toprule
From & To & Setting & Note \\
\midrule
Bench 5 V (or Pi USB) & PPK2 VIN & 5 V & The supply the instrument draws from \\
PPK2 VOUT & ESP32 3V3 pin & source mode, 3300 mV & The DUT's only power source \\
PPK2 GND & DUT GND & common & Also common with the Pi \\
DUT GND & Pi GND & common & One ground reference for the whole bench \\
PPK2 USB & Pi USB-A & CDC & The sample stream and the instrument control \\
ESP32 USB & nothing & disconnected & This is the whole experiment \\
POW-BB & bench rails & labelled & Rails and a common ground, as in P15 \\
Sample rate & PPK2 to Pi & 100 kS/s & Two seconds per firmware state \\
\bottomrule
\end{tabular}
\caption{Wiring. Two cables carry everything: one power path through the instrument, one
USB path to the host. The third cable, the DUT's own USB, stays out of the bench.}
\end{table}

\begin{asciiart}
PPK2 VIN   <--  Pi USB or bench 5V
PPK2 GND   ---  DUT GND --- Pi GND (common)
PPK2 VOUT  -->  ESP32 3V3 pin     source mode, set 3300 mV
PPK2 USB   -->  Pi USB-A
ESP32 USB  DISCONNECTED           this is the whole experiment
\end{asciiart}

Five lines of wiring, and the fifth one is the experiment. The first four put the
instrument in series with the only supply the DUT has, and the fifth removes the supply it
would otherwise prefer. A bench that gets the first four right and the fifth wrong produces
traces that look like measurements, which is worse than producing nothing.

\begin{rulebox}{Two limits that do not move}
The PPK2 in source mode is 5 V-class and about 1 A. Never source a Cat-4 power amplifier
from it: an LTE transmit peak can exceed that current, and the instrument is not a bench
supply for a radio. And never leave the DUT USB cable connected during a capture, because
a second power path turns every number in the run into an unknown fraction of the truth.
\end{rulebox}

\subsection*{Bench layout}

\diagram{p10_bench}{Bench layout. The ESP32 sits on its own with two flying leads and no
USB cable. The PPK2 is between the supply and the DUT, and the only cable leaving the bench
goes to the Pi.}

Lay the bench out so that the disconnected USB cable is visibly elsewhere, not coiled
beside the DUT where somebody will plug it back in between runs. Bring 3V3 and ground to
the ESP32 with two labelled jumpers from the PPK2 output, use the POW-BB discipline from
P15 for any other rail on the bench, and check the ground continuity with a meter before
the first capture rather than after the third.

\begin{note}{The rails come from P15}
This bench uses the POW-BB the way P15 teaches it: labelled rails, a measured 5 V, a
measured 3.3 V and a common ground written in the lab book. Nothing here needs a 5 V pull-up
and nothing here goes onto a GPIO, so the only electrical discipline this lab adds to P15 is
that the DUT has one supply and the instrument is it.
\end{note}

\subsection*{Software design (UML)}

\diagram{p10_uml}{One capture run, across the three firmware states. Each state is flashed,
allowed to settle, captured for two seconds at 100 kS/s, and reduced to a median and a p95.
The run ends with a comparison: if the three traces overlay, the DUT still has a second
power source.}

The activity is deliberately repetitive, because the value of the lab comes from doing the
same thing three times. Flash, settle, capture two seconds, reduce to two numbers, write
them down. The reduction matters as much as the capture: a median says what the state
normally costs, and a p95 says what it costs when the radio is busy, and neither of them is
a screenshot of a trace.

\subsection*{Data flow (ASCII)}

\begin{asciiart}
  bench 5V                 nRF PPK2                      Pi 4
  +---------+          +------------------+          +--------------------+
  |   5 V   |---VIN--->| source mode      |          | nRF Util /         |
  +---------+          |   set 3300 mV    |---USB--->|   Power Profiler   |
                       | 100 kS/s ammeter |   CDC    |   or ppk2-python   |
                       +------------------+          |        |           |
                              |    ^                 |        v           |
                            VOUT   | GND             | 2 s per state      |
                              |    |                 | median and p95     |
                              v    |                 +--------------------+
                       +------------------+
                       | ESP32 NodeMCU    |      USB port:  NOT CONNECTED
                       |   3V3 pin, GND   |      (the whole experiment)
                       | idle | BLE | deep|
                       +------------------+

  One power path, one host path, one ground.  A second power path makes
  every number in the run an unknown fraction of the truth.
\end{asciiart}

\subsection*{Steps}

\begin{steps}
\item \textbf{Install the host tools.} On the Pi, install Nordic's nRF Util with the Power
Profiler command line, or capture CDC samples with the published ppk2-python helpers.
Either path is acceptable; what matters is that the samples land in a file you can reduce
afterwards.

\item \textbf{Wire the rail, then check it.} Bench 5 V into VIN, VOUT to the ESP32 3V3 pin,
grounds common across the PPK2, the DUT and the Pi. Set source mode and 3300 mV before the
DUT is connected, not after.

\item \textbf{Disconnect the DUT USB and keep it disconnected.} The ESP32's own USB cable
stays out of the bench for the whole session. This is the whole experiment.

\item \textbf{Write three tiny sketches.} Radio-off idle, BLE advertise, and deep-sleep. Do
not measure the Arduino blink sketch and call it deep-sleep. Each sketch should do one
thing, so that the state under measurement is the state you named.

\item \textbf{Capture two seconds per state at 100 kS/s.} Let the state settle before the
capture window opens, then take the full two seconds without touching the bench.

\item \textbf{Reduce each capture to two numbers.} Record the median and the p95 for every
state, in a table with the three state names down the side. Quote those numbers afterwards,
not one screenshot of a trace.

\item \textbf{Repeat one state at the end of the run.} Go back to the first state and
capture it again. If the median has moved, something on the bench moved with it: a lead, a
ground, or the state itself. A run that cannot reproduce its own first number is not
finished, whatever the other two states reported.

\item \textbf{Compare the three states before you put the bench away.} Deep-sleep, idle and
BLE advertise should be three visibly different numbers. If they are not, go back to the
DUT USB cable before you go anywhere else.

\item \textbf{Write the honest sentence in the lab book.} Deep-sleep on a dev-board is tens
of microamps to low milliamps because the USB-UART silicon leaks. Write that sentence down
beside the number, so the figure is never quoted later as a bare-module deep-sleep current.
\end{steps}

\subsection*{Acceptance test}

\begin{itemize}
\item Deep-sleep measures tens of microamps to low milliamps on a dev-board, and the lab
book records why: the USB-UART silicon leaks.
\item Active radio measures tens of milliamps.
\item Each state has a median and a p95 written down, not a screenshot.
\item If the three traces overlay, the DUT is still USB-powered. That is a fail, and the
cause is the cable, not the instrument.
\end{itemize}

\subsection*{Practices}

Quote median and p95, not one screenshot. Never source SIM7600 transmit current from the
PPK2. The instrument is 5 V-class and about 1 A in source mode, which covers an ESP32 in
every state this lab measures and does not cover a Cat-4 power amplifier.

Keep this lab free of telemetry. Nothing publishes, nothing subscribes, and no broker is
involved: the numbers are the product. When another lab wants a power figure, it comes here
for it rather than measuring in passing, and it brings its own DUT wiring, because the one
rule that made these numbers valid is that the DUT had exactly one source of power.

Write the bench down as well as the numbers. Which supply fed VIN, what the source voltage
was set to, which sketch each state ran, and how long the state settled before the window
opened. Those four facts are what let somebody repeat the run next month and get the same
table, and they cost four lines in the lab book.

\subsection*{Pitfalls}

\begin{itemize}
\item \textbf{The DUT USB cable still connected.} The symptom is three traces that overlay
at a plausible current. The measurement is of the PPK2's share of a shared rail, which is
not a quantity anyone wants.
\item \textbf{Sourcing a modem from the PPK2.} A Cat-4 transmit peak can exceed about 1 A.
The symptom is a capture with a collapsed rail and a module that resets during transmission.
\item \textbf{Measuring the blink sketch as deep-sleep.} The symptom is a deep-sleep figure
in the milliamps that nobody can reproduce on a bare module.
\item \textbf{Two ground references.} Without a common ground the instrument measures a
loop rather than the DUT. The symptom is a reading that changes when you move a lead.
\item \textbf{Quoting a screenshot.} A trace picture hides the distribution. The symptom is
a power budget built from peaks that were never typical or from averages that were never
measured.
\item \textbf{Capturing before the state settles.} Boot current is not idle current. The
symptom is a median that drifts between repeats of the same state.
\item \textbf{Reporting a dev-board figure as a module figure.} The USB-UART silicon leaks
and the regulator has a quiescent current. Write the sentence in the lab book so the number
is never quoted without it.
\item \textbf{Setting the source voltage with the DUT attached.} Set source mode and
3300 mV first. The symptom is an ESP32 that sees whatever the instrument was last left at.
\item \textbf{Leaving the DUT USB cable coiled beside the board.} It goes back in between
runs, usually without anyone noticing. The symptom is one state in the table that does not
match the others and cannot be reproduced.
\item \textbf{Treating this bench as a telemetry lab.} There is no broker here and no
uplink. The symptom is a capture taken while the DUT is doing something the three sketches
never described, which measures a state nobody named.
\end{itemize}

\subsection*{Sources}

\begin{itemize}
\item Nordic Semiconductor Power Profiler Kit II product page and user guide,
\url{https://www.nordicsemi.com/Products/Development-hardware/Power-Profiler-Kit-2}
\item Nordic nRF Util, with the Power Profiler command line,
\url{https://www.nordicsemi.com/Products/Development-tools/nRF-Util}
\item ppk2-python, the published host helpers for the PPK2 CDC stream,
\url{https://github.com/IRNAS/ppk2-api-python}
\item Espressif ESP32 technical reference and the sleep-mode documentation,
\url{https://docs.espressif.com/projects/esp-idf/en/latest/esp32/api-reference/system/sleep_modes.html}
\item JOY-iT SBC-NodeMCU-ESP32 board documentation,
\url{https://joy-it.net/en/products/SBC-NodeMCU-ESP32}
\item JOY-iT SBC-POW-BB power rail board documentation,
\url{https://joy-it.net/en/products/SBC-POW-BB}
\end{itemize}
